\section{Struktur Dokumen HTML5}

Setiap dokumen HTML5 mengikuti struktur standar yang harus dipahami oleh pengembang web agar browser dapat menginterpretasikan dan menampilkan konten dengan benar. Struktur ini dimulai dengan deklarasi doctype yang memberitahu browser versi HTML yang digunakan. Untuk HTML5, deklarasi ini sangat sederhana: \code{<!DOCTYPE html>}. Deklarasi ini bersifat case-insensitive namun konvensi umum menuliskannya dengan huruf besar. Tanpa deklarasi doctype yang benar, browser mungkin akan beralih ke \textit{quirks mode} yang dapat menyebabkan inkonsistensi dalam rendering CSS dan perilaku JavaScript \cite{whatwgHtml}.

Setelah deklarasi doctype, elemen root \code{<html>} membungkus seluruh konten dokumen. Atribut \code{lang} pada elemen \code{<html>} sangat penting untuk aksesibilitas dan SEO karena memberitahu browser dan mesin pencari bahasa utama konten halaman. Untuk konten berbahasa Indonesia, gunakan \code{<html lang="id">}. Di dalam elemen \code{<html>}, dokumen terbagi menjadi dua bagian utama: \code{<head>} dan \code{<body>}. Bagian \code{<head>} berisi metadata dokumen - informasi tentang halaman itu sendiri yang tidak ditampilkan langsung kepada pengguna, seperti judul, karakter encoding, deskripsi, dan link ke resource eksternal. Bagian \code{<body>} berisi semua konten yang akan terlihat oleh pengguna: teks, gambar, link, form, dan elemen interaktif lainnya \cite{mdnHtmlIntro}.

Struktur hierarki dokumen HTML dapat divisualisasikan sebagai pohon (tree) dengan elemen \code{<html>} sebagai akar (root). Setiap elemen di dalamnya menjadi cabang dengan anak-anak (children) yang lebih spesifik. Memahami struktur pohon ini sangat penting karena CSS selectors dan JavaScript DOM manipulation bekerja berdasarkan hubungan hierarki ini - konsep parent, child, sibling, dan ancestor. Semantik yang benar dalam struktur dokumen juga membantu mesin pencari memahami hierarki informasi dan meningkatkan aksesibilitas bagi pengguna yang mengandalkan teknologi bantu seperti pembaca layar.

\begin{figure}[htbp]
\centering
\begin{tikzpicture}[
  node distance=1.2cm,
  every node/.style={font=\small},
  box/.style={rectangle, draw, minimum width=2.5cm, minimum height=0.6cm, align=center, fill=gray!10},
  child box/.style={rectangle, draw, minimum width=2.2cm, minimum height=0.5cm, align=center, fill=white}
]

\node[box] (html) {\code{<html>} \\ Root};

\node[box, below left=1.5cm and 0.5cm of html] (head) {\code{<head>} \\ Metadata};
\node[box, below right=1.5cm and 0.5cm of html] (body) {\code{<body>} \\ Konten};

\node[child box, below=0.8cm of head] (title) {\code{<title>}};
\node[child box, left=0.1cm of title] (meta) {\code{<meta>}};
\node[child box, right=0.1cm of title] (link) {\code{<link>}};

\node[child box, below=0.8cm of body] (h1) {\code{<h1>}};
\node[child box, left=0.1cm of h1] (nav) {\code{<nav>}};
\node[child box, right=0.1cm of h1] (main) {\code{<main>}};
\node[child box, below=0.5cm of main] (p) {\code{<p>}};
\node[child box, below=0.5cm of nav] (ul) {\code{<ul>} \\ ...};

\draw[->] (html) -- (head);
\draw[->] (html) -- (body);
\draw[->] (head) -- (meta);
\draw[->] (head) -- (title);
\draw[->] (head) -- (link);
\draw[->] (body) -- (nav);
\draw[->] (body) -- (h1);
\draw[->] (body) -- (main);
\draw[->] (main) -- (p);
\draw[->] (nav) -- (ul);

\end{tikzpicture}
\caption{Struktur Hierarki Pohon Dokumen HTML5}
\label{fig:html-tree}
\end{figure}

\begin{webcode}[language=HTML, caption={Struktur Dokumen HTML5 Minimal yang Valid}]
<!DOCTYPE html>
<html lang="id">
  <head>
    <meta charset="UTF-8" />
    <meta name="viewport" content="width=device-width, initial-scale=1.0" />
    <meta name="description" content="Deskripsi halaman untuk SEO" />
    <title>Judul Halaman - Nama Situs</title>
    <link rel="stylesheet" href="styles.css" />
    <link rel="icon" type="image/x-icon" href="favicon.ico" />
  </head>
  <body>
    <header>
      <h1>Selamat Datang</h1>
    </header>
    <main>
      <p>Ini adalah konten utama halaman.</p>
    </main>
    <footer>
      <p>&copy; 2026 Nama Situs</p>
    </footer>
  </body>
</html>
\end{webcode}

\begin{konsep}
\textbf{Komponen Wajib dalam Dokumen HTML5:}
\begin{enumerate}[leftmargin=*]
  \item \textbf{DOCTYPE:} \code{<!DOCTYPE html>} - Deklarasi versi HTML.
  \item \textbf{HTML Root:} \code{<html lang="id">} - Elemen root dengan bahasa.
  \item \textbf{Head:} \code{<head>} - Berisi metadata dan resource.
  \item \textbf{Charset:} \code{<meta charset="UTF-8">} - Encoding karakter.
  \item \textbf{Viewport:} \code{<meta name="viewport">} - Responsif di mobile.
  \item \textbf{Title:} \code{<title>} - Judul yang muncul di tab browser.
  \item \textbf{Body:} \code{<body>} - Berisi semua konten yang terlihat.
\end{enumerate}
\end{konsep}

